\documentclass[a4paper,12pt]{article}
\usepackage[utf8]{inputenc}
\usepackage[T1]{fontenc}
\usepackage{lmodern}
\usepackage{amsmath}
\usepackage{geometry}
\usepackage{booktabs}
\usepackage{graphicx}
\usepackage{enumitem}
\usepackage{xcolor}
\geometry{margin=2.5cm}
\setlength{\parindent}{0pt}
\setlength{\parskip}{0.6em}

\title{IoT - Challenge 4 Homework - Part 1}
\author{Raul Agolli - Team Leader \\\\ Gabriele Lorenzo Singe}
\date{May 2026}

\begin{document}

\maketitle

\section{Exercise 1}

\subsection{Given Data}

The monitoring system is composed of one PAN coordinator and three vibration monitoring nodes.
Each vibration monitoring node processes a sensing window of duration:

\[
T_{\text{window}} = 10 \ \text{s}
\]

The number of vibration peaks detected in each 10-second window follows a Poisson distribution:

\[
X \sim \text{Poisson}(\lambda)
\]

with average rate:

\[
\lambda = 0.4
\]

The system is operated with IEEE 802.15.4 in beacon-enabled mode, considering only the CFP for data transmission.
The CAP is assumed to be zero in this homework:

\[
T_{\text{CAP}} = 0
\]



The nominal bit rate is:

\[
R = 250 \ \text{kbps} = 250000 \ \text{bit/s}
\]

The packet size is:

\[
L = 128 \ \text{bytes}
\]

Since:

\[
1 \ \text{byte} = 8 \ \text{bits}
\]

we have:

\[
L = 128 \cdot 8 = 1024 \ \text{bits}
\]

The payload generated by each node depends on the number of detected vibration peaks:

\[
\begin{cases}
512 \ \text{bytes}, & X=0 \\
2 \ \text{KB}, & X=1 \\
4 \ \text{KB}, & X>1
\end{cases}
\]

Assuming:

\[
1 \ \text{KB} = 1024 \ \text{bytes}
\]


Since one packet has size \(128\) bytes, the number of packets generated in a 10-second window in the three cases is:

\[
n_0 = \frac{512}{128} = 4
\]

\[
n_1 = \frac{2048}{128} = 16
\]

\[
n_2 = \frac{4096}{128} = 32
\]

where \(n_0\), \(n_1\), and \(n_2\) correspond respectively to the cases \(X=0\), \(X=1\), and \(X>1\).

\subsection{Output Rate Probability Mass Function}

For a Poisson random variable, the probability of observing \(k\) events is:

\[
P(X=k)=\frac{\lambda^k e^{-\lambda}}{k!}
\]

Therefore, with:

\[
\lambda = 0.4
\]

we obtain:

\[
P(X=0)=\frac{0.4^0 e^{-0.4}}{0!}=0.6703
\]

For one vibration peak:

\[
P(X=1)=\frac{0.4^1 e^{-0.4}}{1!}
=0.2681\]

For more than one vibration peak:

\[
P(X>1)=1-P(X=0)-P(X=1)=0.0616
\]

The corresponding output rates, expressed in packets per second over the 10-second sensing window, are:

\[
r_0 = \frac{4}{10} = 0.4 \ \text{packets/s}
\]

\[
r_1 = \frac{16}{10} = 1.6 \ \text{packets/s}
\]

\[
r_2 = \frac{32}{10} = 3.2 \ \text{packets/s}
\]

Therefore, the probability mass function of the output rate is:

\[
\boxed{P(r=r_0)=0.6703}\\
\boxed{P(r=r_1)=0.2681}\\
\boxed{P(r=r_2)=0.0616}
\]

\subsection{CFP Slot Assignment}

Since the system is operated in beacon-enabled mode and the CFP is used for scheduled transmissions, the CFP slot assignment is dimensioned considering the worst-case output rate.

In the worst case, a vibration monitoring node detects more than one vibration peak:

\[
X>1
\]

and therefore generates a payload of:

\[
4 \ \text{KB} = 4096 \ \text{bytes}
\]

This corresponds to:

\[
32 \ \text{packets in } 10 \ \text{s}
\]

Thus, the worst-case packet generation rate is:

\[
r_2 = \frac{32}{10} = 3.2 \ \text{packets/s}
\]

Since the beacon interval is:

\[
BI = 2.5 \ \text{s}
\]

the number of packets generated by one node during one beacon interval is:

\[
N_{\text{CFP slots/node}} = r_2 \cdot BI = 8
\]

Since one packet fits exactly in one CFP slot, each node requires:

\[
\boxed{N_{\text{CFP slots/node}} = 8}
\]

Since the system has three vibration monitoring nodes, the total number of CFP slots is:

\[
N_{\text{CFP slots}} = 3 \cdot 8 = 24
\]

Therefore:

\[
\boxed{N_{\text{CFP slots}} = 24}
\]

A consistent CFP slot assignment is 8 CFP slots per node

\subsection{Slot Duration}

The duration of one CFP slot is equal to the time needed to transmit one packet.

The packet size is:

\[
L = 128 \ \text{bytes} = 1024 \ \text{bits}
\]

The nominal bit rate is:

\[
R = 250000 \ \text{bit/s}
\]

Therefore, the slot duration is:

\[
T_s = \frac{L}{R} = 4.096 \ \text{ms}
\]


Therefore:

\[
\boxed{T_s = 4.096 \ \text{ms}}
\]

\subsection{Active Time, Inactive Time and Duty Cycle}

The CFP contains:

\[
N_{\text{CFP slots}} = 24 \ slots
\]


However, the active part of the beacon-enabled superframe also includes the initial beacon signal.
Since the CAP is zero in this homework, the active part is composed of:

\[
\text{Beacon} + \text{CFP}
\]

Therefore, the total number of active slots is:

\[
N_{\text{active slots}} = 1 + N_{\text{CFP slots}} = 25
\]

The active time is:

\[
T_{\text{active}} = N_{\text{active slots}} \cdot T_s = 102.4 \ \text{ms}
\]

Therefore:

\[
\boxed{T_{\text{active}} = 102.4 \ \text{ms}}
\]

The beacon interval is:

\[
BI = 2.5 \ \text{s} = 2500 \ \text{ms}
\]

Thus, the inactive time is:

\[
T_{\text{inactive}} = BI - T_{\text{active}} = 2397.6 \ \text{ms}   = 2.3976 \ \text{s}
\]

Therefore:

\[
\boxed{T_{\text{inactive}} = 2.3976 \ \text{s}}
\]

The duty cycle is defined as:

\[
DC = \frac{T_{\text{active}}}{BI} = 0.04096
\]

Expressed as a percentage:

\[
DC = 4.096\%
\]

Therefore:

\[
\boxed{DC = 4.096\%}
\]

\subsection{Maximum Number of Additional Nodes}

The system must keep the duty cycle below:

\[
10\%
\]

Since the beacon interval is:

\[
BI = 2.5 \ \text{s} = 2500 \ \text{ms}
\]

the maximum allowed active time is:

\[
T_{\text{active,max}} = 0.10 \cdot 2500= 250 \ \text{ms}
\]


The current active time is:

\[
T_{\text{active}} = 102.4 \ \text{ms}
\]

Therefore, the remaining active time available before reaching the \(10\%\) duty-cycle limit is:

\[
T_{\text{remaining}} = T_{\text{active,max}} - T_{\text{active}}= 147.6 \ \text{ms}
\]

Each additional vibration monitoring node requires 8 CFP slots where each slot lasts 4.096 ms.


Therefore, the active time required by one additional node is:

\[
T_{\text{node}} = 8 \cdot 4.096= 32.768 \ \text{ms}
\]


The maximum number of additional nodes is:

\[
\Delta M =
\left\lfloor
\frac{T_{\text{remaining}}}{T_{\text{node}}}
\right\rfloor
=
\left\lfloor
4.50
\right\rfloor
=4
\]


Therefore:

\[
\boxed{\Delta M = 4}
\]

\subsection{Final Results}

Using worst-case CFP dimensioning, a beacon interval of \(BI=2.5\,\text{s}\), and including the initial beacon signal in the active time, the final results are:

\[
\textbf{Probabilities}
\qquad
\boxed{
\begin{aligned}
P(r=r_0) &= 0.6703 \\
P(r=r_1) &= 0.2681 \\
P(r=r_2) &= 0.0616
\end{aligned}
}
\]

\[
\textbf{CFP slots}
\qquad
\boxed{
\begin{aligned}
N_{\text{CFP slots/node}} &= 8 \\
N_{\text{CFP slots}} &= 24
\end{aligned}
}
\]

\[
\textbf{Timing and Duty Cycle}
\qquad
\boxed{
\begin{aligned}
T_s &= 4.096 \ \text{ms} \\
T_{\text{active}} &= 102.4 \ \text{ms} \\
T_{\text{inactive}} &= 2.3976 \ \text{s} \\
DC &= 4.096\%
\end{aligned}
}
\]

\[
\textbf{Additional Nodes}
\qquad
\boxed{
\Delta M = 4
}
\]
\end{document}